\section{A finite-endpoint Mellin--Hurwitz transform}
\label{ag:sec:hurwitz}

Fix $r\ge2$ and initially take real $a>b>0$. Define
\begin{equation}\label{ag:eq:Hmellin}
 \HH_r(s;a,b)=\frac1{\Gamma(s)}\int_0^\infty
       \frac{t^{s-1}e^{-at}g_r(bt)}{1-e^{-t}}\dd t,
                  \qquad \Re s>1.
\end{equation}
At zero the integrand is $O(t^{\Re s-2})$, and at infinity the
exponential decay is governed by $a-b>0$. These conditions justify
ordinary integration and all fixed-order spectral derivatives in the
stated half-plane.

The Hurwitz zeta convention is
$\zeta(s,x)=\sum_{n\ge0}(n+x)^{-s}$ for $\Re s>1$, continued
meromorphically in $s$. Its standard Mellin representation and spatial
derivative are recorded in \cite{DLMFzeta}. They also follow directly
from the series in their common domain.

\begin{theorem}[The beta average and its finite endpoints]
\label{ag:thm:hurwitz}
With $c_r$ as in \eqref{ag:eq:beta-kernel},
\begin{equation}\label{ag:eq:Hcompact}
 \HH_r(s;a,b)=\frac{c_r}{2}\int_{-1}^1
                    (1-v^2)^r\zeta(s,a+bv)\dd v.
\end{equation}
This provides meromorphic continuation to every $s$, with the single
possible pole $s=1$, whose residue is
\begin{equation}\label{ag:eq:Hresidue}
 g_r(0)=\frac{(-1)^r}{(2r+2)!}.
\end{equation}
Define
\begin{equation}\label{ag:eq:weights}
 w_{rj}(b)=\frac{(-1)^j j!\binom r{j-r}}
 {2^{j+2}r!(r+1)!b^{j+1}},\qquad r\le j\le2r.
\end{equation}
Then, as an identity of meromorphic functions,
\begin{equation}\label{ag:eq:Hendpoints}
 \boxed{
 \HH_r(s;a,b)=\sum_{j=r}^{2r}w_{rj}(b)
 \frac{\zeta(s-j-1,a+b)-(-1)^j\zeta(s-j-1,a-b)}{(1-s)_{j+1}}.}
\end{equation}
At apparent poles of the separate summands, the entire sum is to be
continued jointly. No pole other than $s=1$ survives.
\end{theorem}
\begin{proof}
Insert \eqref{ag:eq:beta-kernel} into \eqref{ag:eq:Hmellin}.
Absolute convergence permits interchanging the integrals. Combining
the two exponentials in the hyperbolic cosine yields
\eqref{ag:eq:Hcompact}. On every compact subset of $\Re x>0$,
$\zeta(s,x)-1/(s-1)$ is holomorphic and locally bounded in $s$.
Its compact $v$-integral is therefore entire, proving the continuation
and the residue. The integral of the weight is the one already used
in \eqref{ag:eq:kernel-coeff}.

For the endpoint identity let $P(v)=(1-v^2)^r$. Repeated spatial
antidifferentiation gives
\[
 \frac{\mathrm d^{j+1}}{\mathrm dv^{j+1}}
 \frac{\zeta(s-j-1,a+bv)}{b^{j+1}(1-s)_{j+1}}
 =\zeta(s,a+bv),
\]
initially for $s$ outside the finite exceptional set. Integrating by
parts until the polynomial vanishes leaves
\[
 \sum_{j=0}^{2r}\frac{(-1)^j}{b^{j+1}(1-s)_{j+1}}
       [P^{(j)}(v)\zeta(s-j-1,a+bv)]_{v=-1}^{v=1}.
\]
The derivatives with $j<r$ vanish at both endpoints, and for $r\le j\le2r$,
\begin{equation}\label{ag:eq:weight-derivative}
 P^{(j)}(1)=(-1)^r j!\binom r{j-r}2^{2r-j},\qquad
 P^{(j)}(-1)=(-1)^jP^{(j)}(1).
\end{equation}
Multiplication by $c_r/2$ proves \eqref{ag:eq:Hendpoints}. Analytic
continuation establishes it at every remaining regular point.
\end{proof}

The formulas extend holomorphically in $a,b$ when $\Re a>|b|$, using
the right-half-plane Hurwitz branch. They are even in $b$. At $b=0$
the compact formula gives the confluent value
\begin{equation}\label{ag:eq:Hcollision}
 \HH_r(s;a,0)=g_r(0)\zeta(s,a).
\end{equation}
The negative powers of $b$ in \eqref{ag:eq:Hendpoints} must cancel;
one should not substitute $b=0$ separately in those terms.

\subsection{All logarithmically weighted Mellin integrals}
Let $\mathsf Y_m$ be the complete exponential Bell polynomial, defined
by $\exp(\sum_{j\ge1}x_jz^j/j!)=\sum_{m\ge0}\mathsf Y_m(x_1,\ldots,x_m)z^m/m!$.
Differentiating the normalized Mellin factor yields, for $\Re s>1$,
\begin{align}
 \partial_s^m\HH_r(s;a,b)
 =\frac1{\Gamma(s)}\int_0^\infty&
 \frac{t^{s-1}e^{-at}g_r(bt)}{1-e^{-t}}\notag\\[-2pt]
 &\quad\times\mathsf Y_m\bigl(\log t-\psi(s),-\psi'(s),
                      \ldots,-\psi^{(m-1)}(s)\bigr)\dd t.
 \label{ag:eq:log-Mellin}
\end{align}
For $m=0$ the Bell factor is one. For $m=1,2$ the factors are
$\log t-\psi(s)$ and $(\log t-\psi(s))^2-\psi'(s)$.
This is an exact normalization, not the unnormalized rule of merely
inserting $(\log t)^m$. Endpoint differentiation of
\eqref{ag:eq:Hendpoints} therefore gives finite evaluations of all
these convergent logarithmic integrals.
